\chapter{Pap Smear (Cervical Cancer Screening)}

\section*{What Is This Test?}
A Pap smear, or cervical cytology test, examines cells collected from the cervix for precancerous changes and cervical cancer. It may be performed alone or together with a high-risk human papillomavirus (HPV) test, depending on age, risk, and local screening guidance. The Pap test is a screening test, not a test for every cause of vaginal symptoms.

\section*{Alternative Names}
Pap Test; Pap Smear; Cervical Cytology; Cervical Screening; Liquid-Based Cytology; Cervical Cancer Screening

\section*{Why It Is Ordered}
Routine cervical-cancer screening at guideline-recommended intervals, follow-up of an abnormal HPV or Pap result, and surveillance after treatment of cervical precancer. Screening schedules vary by age, prior results, immune status, and history of cervical disease.

\section*{How This Test Is Performed}
During a pelvic examination, a speculum gently separates the vaginal walls and a clinician uses a soft brush or spatula to collect cells from the cervix. The sample is placed in a liquid vial or on a slide and examined by a cytology laboratory. HPV testing may be performed from the same sample.

\section*{Preparation, Risks, and Limitations}
Follow local instructions; some clinics advise avoiding vaginal medicines, douching, and sexual intercourse for a short period before collection. Mild discomfort or light spotting may occur. An abnormal result does not mean cancer is present; it may reflect HPV infection, inflammation, or a sampling issue and usually requires risk-based follow-up. A normal Pap result does not eliminate the need for future screening.

\section*{References}
\begin{enumerate}
  \item MedlinePlus. Pap Smear.
  \item United States Preventive Services Task Force. Cervical Cancer: Screening.
\end{enumerate}

\clearpage
\section*{Understanding Results and Follow-up}
The laboratory report should identify the specimen, method, units, reference interval or decision limit, and whether the result is positive, negative, abnormal, or inconclusive. Reference intervals vary with age, sex, pregnancy, specimen type, assay, and laboratory; a result outside a range is not automatically a diagnosis, and a result within a range does not always exclude disease. Discuss unexpected results with the ordering clinician, who can decide whether repeat or confirmatory testing, treatment, or referral is appropriate. Do not change medicines or delay urgent care based on an isolated result.


\section*{Regulatory Status and Authorization Date}
Regulatory status applies to the specific assay, instrument, manufacturer, and intended use rather than to the test name alone. No single approval date can be assigned to this test category. For a commercial assay, verify the current product in the relevant regulator's database: the U.S. Food and Drug Administration (FDA) clearance, approval, or authorization; India's Central Drugs Standard Control Organisation (CDSCO) licence or permission; the European Union In Vitro Diagnostic Regulation (EU IVDR) CE certificate issued through the applicable conformity-assessment route; or the United Kingdom Medicines and Healthcare products Regulatory Agency (MHRA) registration and UKCA/CE status. Laboratory-developed tests may not have an individual FDA, CDSCO, EU IVDR, or MHRA approval date.
\clearpage
